% SGA 3, Expose IV, English reconstruction.
% Scope: all of Section 5, stopping immediately before the Section 6 heading.
% Authority: Polo--Gille Exp4-13oct24.pdf, component-local pages 55--61
% (combined-reader pages 243--249).
% Jacob Reinhold's CC BY 4.0 Markdown lineage was used only as a phrasing
% comparison. The Polo--Gille PDF controls prose, formulas, numbering, and notes.

\setcounter{footnote}{47}

% Source page: Exp4 component-local p. 55; combined-reader p. 243.
\SGARefTarget{sga3:iv:section:5}\section*{5. 取商与代数结构}
\addcontentsline{toc}{section}{5. 取商与代数结构}

\SGARefTarget{sga3:iv:subsection:5:1}\subsection*{5.1. 主齐性丛}
\addcontentsline{toc}{subsection}{\SGARefLink{sga3:iv:subsection:5:1}{5.1}. 主齐性丛}

\noindent\SGARefTarget{sga3:iv:definition:5:1:0}\textbf{定义 5.1.0.}\footnote{编者注: 引入编号\SGARefLink{sga3:iv:definition:5:1:0}{5.1.0}, 以便后文引用.}
回顾 (III \SGARefLink{sga3:iii:definition:0:1}{0.1}), 带有
\(H\) 的右作用的对象 \(X\) 称为 \(H\) 下的\emph{形式主齐性对象 (formally principal homogeneous)}\footnote{编者注: 也称为``伪挠子 (pseudo-torsor)''; 参见 EGA IV\textsubscript{4}, 16.5.15.
更一般的形式齐性对象 (未必主齐性) 的概念, 定义于本篇末尾的补遗
\SGARefLink{sga3:iv:definition:6:7:1}{6.7.1}.} 若如下典范函子态射
\[
  X\times H\longrightarrow X\times X,
  \qquad (x,h)\longmapsto(x,xh),
\]
为同构. 等价地 (参见\emph{同上}), 对每个
\(S\in\operatorname{Ob}\mathscr C\), 集合 \(X(S)\) 在 \(H(S)\) 下形式主齐性:
它或为空集, 或为 \(H(S)\) 下的主齐性集.
特别地, 当 \(H\) 通过右平移作用于自身时, 它在自身作用下形式主齐性.

\medskip
\noindent\SGARefTarget{sga3:iv:definition:5:1:1}\textbf{定义 5.1.1.}
带有 \(H\) 作用的对象 \(X\) 称为\emph{平凡的 (trivial)},
若它作为 \(H\)-对象同构于带有平移作用的 \(H\).

\medskip
\noindent\SGARefTarget{sga3:iv:proposition:5:1:2}\textbf{命题 5.1.2.}
\textit{若 \(X\) 在 \(H\) 下形式主齐性, 则存在
\(\Gamma(H)\) 下的主齐性集之间的同构}
\[
  \Gamma(X)\xrightarrow{\sim}
  \operatorname{Isom}_{H\text{-obj.}}(H,X)
\]
\textit{}

对 \(X\) 的截面 \(x\), 关联由集合上的映射 \(h\mapsto xh\) 定义的态射
\(H\to X\). 归结为集合论情形, 立即得到\SGARefLink{sga3:iv:proposition:5:1:2}{此断言}.

\medskip
\noindent\SGARefTarget{sga3:iv:corollary:5:1:3}\textbf{推论 5.1.3.}
\textit{存在 \(H\)-对象的同构}
\[
  X\xrightarrow{\sim}\operatorname{Isom}_{H\text{-obj.}}(H,X).
\]

\medskip
\noindent\SGARefTarget{sga3:iv:corollary:5:1:4}\textbf{推论 5.1.4.}
\textit{带有群作用的对象平凡, 当且仅当它形式主齐性且有截面.}

% Source page: Exp4 component-local p. 56; combined-reader p. 244.
\medskip
\noindent\SGARefTarget{sga3:iv:definition:5:1:5}\textbf{定义 5.1.5.}
设 \(\mathscr C\) 为配备拓扑的范畴.
带有 \(S\)-群 \(H\) 作用的 \(S\)-对象 \(X\) 称为
\emph{\(H\) 下的主齐性丛}, 若它局部平凡, 即满足下述等价条件:\footnote{编者注: 此时也称 \(X\) 为 \(H\)-挠子 (torsor).}
\begin{enumerate}
  \item[(i)]\SGARefTarget{sga3:iv:definition:5:1:5:item:i:01} 所有使函子 \(X\times_ST\) 在 \(H\times_ST\) 下平凡的态射
\(T\to S\) 所组成的族, 是 \(S\) 的一个加细;
  \item[(ii)]\SGARefTarget{sga3:iv:definition:5:1:5:item:ii:01} 对于定义所给拓扑的一个预拓扑, 存在覆盖族 \(\{S_i\to S\}\),
使得对每个 \(i\), \(S_i\)-函子 \(X\times_SS_i\) 在 \(S_i\)-群函子
\(H\times_SS_i\) 下平凡, 等价地, 它在 \(S_i\) 上有截面.
\end{enumerate}


\medskip
\noindent\SGARefTarget{sga3:iv:proposition:5:1:6}\textbf{命题 5.1.6.}
\textit{设 \(\mathscr C\) 为配备拓扑 \(\mathcal T\) 的范畴,
\((M)\) 为满足\SGARefLink{sga3:iv:para:4:6:3}{4.6.3} 的公理
\SGARefLink{sga3:iv:para:4:6:3:item:a:01}{(a)}--\SGARefLink{sga3:iv:para:4:6:3:item:e:mathcal:t:01}{\((e_{\mathcal T})\)}
的 \(\mathscr C\) 中的一族态射.
设 \(H\) 为结构态射 \(H\to S\) 属于 \((M)\) 的 \(S\)-群,
\(P\) 为带有 \(H\) 作用的 \(S\)-对象. 下列条件等价:}
\begin{enumerate}
  \item[(i)]\SGARefTarget{sga3:iv:proposition:5:1:6:item:i:01} \textit{\(P\) 为 \(H\) 下的主齐性丛
(定义\SGARefLink{sga3:iv:definition:5:1:5}{5.1.5});}
  \item[(ii)]\SGARefTarget{sga3:iv:proposition:5:1:6:item:ii:01} \textit{\(P\) 在 \(H\) 下形式主齐性, 且其结构态射 \(P\to S\) 属于 \((M)\);}
  \item[(iii)]\SGARefTarget{sga3:iv:proposition:5:1:6:item:iii:01} \textit{存在属于 \((M)\) 的态射 \(S'\to S\), 使得从 \(S\) 到 \(S'\)
基变换后, \(P\) 成为平凡对象; 即 \(P\times_SS'\) 在 \(H\times_SS'\) 下平凡;}
  \item[(iv)]\SGARefTarget{sga3:iv:proposition:5:1:6:item:iv:01} \textit{\(H\) 自由且 \((M)\)-有效地作用于 \(P\), 且商 \(P/H\) 同构于 \(S\).}
\end{enumerate}

首先, \SGARefLink{sga3:iv:proposition:5:1:6:item:ii:01}{(ii)}
与\SGARefLink{sga3:iv:proposition:5:1:6:item:iv:01}{(iv)} 等价.
两种情形下 \(P\to S\) 都属于 \((M)\), 从而可平方,
这保证纤维积 \(H\times_SP\) 和 \(P\times_SP\) 可表.
条件\SGARefLink{sga3:iv:proposition:5:1:6:item:ii:01}{(ii)} 显然蕴涵\SGARefLink{sga3:iv:proposition:5:1:6:item:iii:01}{(iii)}:
可取 \(S'=P\), 因为形式主齐性与\SGARefLink{sga3:iv:corollary:5:1:4}{5.1.4} 表明,
\(P\times_SP\) 在 \(H\times_SP\) 下平凡, 其对角截面给出所需截面.
条件\SGARefLink{sga3:iv:proposition:5:1:6:item:iii:01}{(iii)} 显然蕴涵\SGARefLink{sga3:iv:proposition:5:1:6:item:i:01}{(i)},
因为由公理\SGARefLink{sga3:iv:para:4:6:3:item:d:mathcal:t:01}{\((d_{\mathcal T})\)},
\(\{S'\to S\}\) 是覆盖族.

只需证明\SGARefLink{sga3:iv:proposition:5:1:6:item:i:01}{(i)}\(\Rightarrow\)\SGARefLink{sga3:iv:proposition:5:1:6:item:ii:01}{(ii)}.
层态射
\[
  P\times_SH\longrightarrow P\times_SP
\]
局部为同构, 从而是同构 (\SGARefLink{sga3:iv:corollary:4:5:8}{4.5.8});
故 \(P\) 形式主齐性. 结构态射 \(P\to S\) 局部同构于属于 \((M)\)
的结构态射 \(H\to S\), 因而由公理
\SGARefLink{sga3:iv:para:4:6:3:item:e:mathcal:t:01}{\((e_{\mathcal T})\)}, 它也属于 \((M)\).

\SGARefLink{sga3:iv:proposition:5:1:6:item:i:01}{(i)} 与\SGARefLink{sga3:iv:proposition:5:1:6:item:iv:01}{(iv)}
的等价性有如下推广.

\medskip
\noindent\SGARefTarget{sga3:iv:proposition:5:1:7}\textbf{命题 5.1.7.}
\textit{对 \(\mathscr C\) 和 \((M)\) 作相同假设,
设 \(H\) 为 \(S\)-群, \(X\) 为带有 \(H\) 的右作用的 \(S\)-对象.
假设 \(H\to S\) 属于 \((M)\). 下列条件等价:}
\begin{enumerate}
  \item[(i)]\SGARefTarget{sga3:iv:proposition:5:1:7:item:i:01} \textit{\(H\) 自由且 \((M)\)-有效地作用于 \(X\);}
  \item[(ii)]\SGARefTarget{sga3:iv:proposition:5:1:7:item:ii:01} \textit{存在与 \(H\) 的作用在 \(X\) 上定义的等价关系相容的
\(S\)-态射 \(p:X\to Y\), 使得 \(H\times_SY\) 在 \(X\) 上诱导的
\(Y\) 上的作用, 将 \(X\) 赋成 \(Y\) 上的 \(H_Y\) 下的主齐性丛.}
\end{enumerate}
\textit{在这些条件下, \(p\) 将 \(Y\) 与商 \(X/H\) 等同.}

% Source page: Exp4 component-local p. 57; combined-reader p. 245.
若 \(p:X\to Y\) 与 \(H\) 的作用相容,
则由同构
\[
  H_Y\times_YX\xrightarrow{\sim}H\times_SX.
\]
可知, \(H\times_SY\) 在 \(X\) 上诱导的 \(Y\) 上的作用,
与 \(H\) 的原作用在 \(X\) 上定义同一个等价关系.
由此观察及前述等价性\SGARefLink{sga3:iv:proposition:5:1:6:item:iv:01}{(iv)}\(\Leftrightarrow\)\SGARefLink{sga3:iv:proposition:5:1:6:item:i:01}{(i)},
即得\SGARefLink{sga3:iv:proposition:5:1:7}{此命题}.

\medskip
\noindent\SGARefTarget{sga3:iv:corollary:5:1:7:1}\textbf{推论 5.1.7.1.}\footnote{编者注: 此特殊情形作为推论加入, 将在后续各篇中多次使用.}
\textit{设 \(\mathscr C\) 为具有终对象和纤维积的范畴,
配备比典范拓扑更粗的拓扑 \(\mathcal T\).
设 \(f:G\to H\) 为 \(\mathscr C\)-群的态射,
令 \(K=\operatorname{Ker}(f)\), 并假设 \(f\) 是 \(\mathcal T\) 的覆盖态射.
则 \(H\) 表出层商 \(G/K\), 且 \(f\) 是 \(K_H\)-挠子.
等价地, 可以说 \(G\) 是 \(H\) 上的 \(K\)-挠子.}

事实上, 由于 \(f\) 是覆盖态射, 它是普遍有效满态射
(\SGARefLink{sga3:iv:proposition:4:4:3}{4.4.3}).
故由\SGARefLink{sga3:iv:scholium:3:3:3:1}{3.3.3.1}, \(H\) 是 \(G\) 对等价关系
\[
  \mathscr R(f)=G\times_HG.
\]
的商. 此外, 态射
\[
  G\times K\longrightarrow G\times_HG,
  \qquad (g,k)\longmapsto(g,gk),
\]
是 \(K_G=G\times_HK_H\)-对象的同构; 其逆为
\((g,g')\mapsto(g,g^{-1}g')\). 因此 \(\mathscr R(f)\) 是 \(K\) 定义的等价关系.
由于 \(f:G\to H\) 是覆盖态射, 由\SGARefLink{sga3:iv:proposition:5:1:6:item:ii:01}{5.1.6(ii)},
或直接由定义\SGARefLink{sga3:iv:definition:5:1:5:item:ii:01}{5.1.5(ii)}, 它是 \(K_H\)-挠子.

现在可以对群作用取商的情形加强定理\SGARefLink{sga3:iv:theorem:4:6:6}{4.6.6}.

\medskip
\noindent\SGARefTarget{sga3:iv:proposition:5:1:8}\textbf{命题 5.1.8.}
\textit{在\SGARefLink{sga3:iv:proposition:5:1:7}{5.1.7} 的假设下,
用 \(F_0\) 表示如下定义的 \(S\) 上的函子: 对每个 \(S'\to S\),
\(F_0(S')\) 是 \(X\times_SS'\) 的那些可表 \(S'\)-子函子 \(Z\) 组成的集合,
这些子函子在 \(H\times_SS'\) 下稳定, 且对诱导的作用,
是此 \(S'\)-群下的主齐性丛 (\SGARefLink{sga3:iv:lemma:3:2:2}{3.2.2}).}
\begin{enumerate}
  \item[(i)]\SGARefTarget{sga3:iv:proposition:5:1:8:item:i:01} \textit{下列条件等价:}
  \begin{enumerate}
    \item[(a)]\SGARefTarget{sga3:iv:proposition:5:1:8:item:a:01} \textit{\(H\) 在 \(X\) 上的作用是 \((M)\)-有效且自由的;}\footnote{编者注: 增加了``且自由''一语.}
    \item[(b)]\SGARefTarget{sga3:iv:proposition:5:1:8:item:b:01} \textit{\(F_0\) 可表.}
  \end{enumerate}
  \textit{在这些条件下, \(F_0=X/H\).}
  \item[(ii)]\SGARefTarget{sga3:iv:proposition:5:1:8:item:ii:01} \textit{设 \((N)\) 为对基变换稳定的态射族, 使得对每个覆盖族
\(\{S_i'\to S'\}\) 和 \((N)\) 中的每个态射族 \(\{T_i\to S_i'\}\),
\(T_i\) 上关于 \(\{S_i'\to S'\}\) 的每个下降资料都是有效的.
假设 \(X\times_SH\to X\times_SX\) 属于 \((N)\), 且 \(X\) 可平方.
则与 \(X\times_SX\) 的子对象 \(X\times_SH\) 对应的元素
\(p\in\operatorname{Hom}(X,F_0)\), 将 \(F_0\) 与层商 \(X/H\) 等同.}
\end{enumerate}

% Source page: Exp4 component-local p. 58; combined-reader p. 246.
\SGARefTarget{sga3:iv:subsection:5:2}\subsection*{5.2. 群结构与取商}
\addcontentsline{toc}{subsection}{\SGARefLink{sga3:iv:subsection:5:2}{5.2}. 群结构与取商}

本小节考察群对其子群的商 (G/H) 上可赋予的代数结构.
先在 \(\mathscr C\) 上关于任意拓扑的层范畴中讨论.
再取典范拓扑并使用\SGARefLink{sga3:iv:lemma:4:5:12}{4.5.12},
即可得到 \(\mathscr C\) 中普遍有效取商的结果.

\medskip
\noindent\SGARefTarget{sga3:iv:proposition:5:2:1}\textbf{命题 5.2.1.}
\textit{设 \(u:H\to G\) 为群层的单态射.
商层 \(G/H\) 上存在唯一的 \(G\)-对象结构, 使得典范态射
\[
  p:G\longrightarrow G/H
\]
为 \(G\)-对象的态射. 此结构关于对 \((G,H)\) 具有函子性:
给定交换图
\[
  \begin{array}{ccc}
    H & \longrightarrow & G \\
    \downarrow && \downarrow\scriptstyle f \\
    H' & \longrightarrow & G',
  \end{array}
\]
\SGARefLink{sga3:iv:lemma:3:2:3}{3.2.3} 的态射 \(f:G/H\to G'/H'\)
与作用群之间的态射 \(f\) 相容.}

事实上, \(G/H\) 是下列预层的伴随层:
\[
  i(G)/i(H):S\longmapsto G(S)/H(S).
\]
由于函子 \(a\) 左正合, 它将带有群作用的对象映到带有群作用的对象.
预层 \(i(G)/i(H)\) 是 \(i(G)\)-对象, 因此
\[
  G/H=a\bigl(i(G)/i(H)\bigr)
\]
是 \(a(i(G))=G\)-对象. 此构造显然具有所述的全部性质.

\medskip
\noindent\SGARefTarget{sga3:iv:corollary:5:2:2}\textbf{推论 5.2.2.}
\textit{设 \(u:H\to G\) 为 \(\mathscr C\)-群的单态射,
并假设 \(H\) 在 \(G\) 上的作用普遍有效.
商对象 \(G/H\in\operatorname{Ob}\mathscr C\) 上存在唯一的 \(G\)-对象结构,
使得 \(p:G\to G/H\) 为带有作用的对象的态射.
在\SGARefLink{sga3:iv:proposition:5:2:1}{前一命题}的意义下,
此结构关于满足 \(H\) 普遍有效地作用于 \(G\) 的对 \((H,G)\) 具有函子性.}

\medskip
\noindent\SGARefTarget{sga3:iv:proposition:5:2:3}\textbf{命题 5.2.3.}
\textit{设 \(u:H\to G\) 为群层的单态射, 将 \(H\) 与 \(G\) 的一个正规子群层等同.
商层 \(G/H\) 上存在唯一的群层结构, 使典范态射 \(p:G\to G/H\) 为群态射.
此结构关于满足 \(H\) 正规的对 \((H,G)\) 具有函子性.}

证明与\SGARefLink{sga3:iv:proposition:5:2:1}{5.2.1} 的证明类似.

\medskip
\noindent\SGARefTarget{sga3:iv:corollary:5:2:4}\textbf{推论 5.2.4.}
\textit{设 \(u:H\to G\) 为 \(\mathscr C\)-群的单态射,
将 \(H\) 与 \(G\) 的一个正规子群等同.
假设 \(H\) 在 \(G\) 上的作用普遍有效.
商对象 \(G/H\in\operatorname{Ob}\mathscr C\) 上存在唯一的群结构,
使典范态射 \(G\to G/H\) 为群态射.
此结构关于满足 \(H\) 正规且普遍有效地作用的对 \((H,G)\) 具有函子性.}

% Source page: Exp4 component-local p. 59; combined-reader p. 247.
可以更明确地刻画 \(G/H\) 的群结构.

\medskip
\noindent\SGARefTarget{sga3:iv:proposition:5:2:5}\textbf{命题 5.2.5.}
\textit{在\SGARefLink{sga3:iv:corollary:5:2:4}{5.2.4} 的假设下,
设 \(K\) 为 \(\mathscr C\)-群, \(f:G\to K\) 为态射. 下列条件等价:}
\begin{enumerate}
  \item[(i)]\SGARefTarget{sga3:iv:proposition:5:2:5:item:i:01} \textit{\(f\) 是与 \(H\) 定义的等价关系相容的群态射;}
  \item[(ii)]\SGARefTarget{sga3:iv:proposition:5:2:5:item:ii:01} \textit{\(f\) 是群态射, 且其限制 \(H\to K\) 为平凡态射;}
  \item[(iii)]\SGARefTarget{sga3:iv:proposition:5:2:5:item:iii:01} \textit{\(f\) 经由群态射 \(G/H\to K\) 分解.}
\end{enumerate}
\textit{特别地, 存在关于群 \(K\) 具有函子性的同构}
\[
  \operatorname{Hom}_{\mathscr C\text{-gr.}}(G/H,K)
  \xrightarrow{\sim}
  \bigl\{f\in\operatorname{Hom}_{\mathscr C\text{-gr.}}(G,K)
  \mid f\circ u=e\bigr\}.
\]

\SGARefLink{sga3:iv:proposition:5:2:5:item:i:01}{(i)} 与\SGARefLink{sga3:iv:proposition:5:2:5:item:ii:01}{(ii)}
的等价性可在集合论层面证明,
而\SGARefLink{sga3:iv:proposition:5:2:5:item:iii:01}{(iii)}\(\Rightarrow\)\SGARefLink{sga3:iv:proposition:5:2:5:item:ii:01}{(ii)} 显然成立.
\SGARefLink{sga3:iv:proposition:5:2:5:item:iii:01}{(iii)} 与\SGARefLink{sga3:iv:proposition:5:2:5:item:ii:01}{(ii)}
的等价性由公式
\[
  \operatorname{Hom}(G/H,K)
  \simeq \operatorname{Hom}\bigl(i(G)/i(H),K\bigr)
\]
及 \(G/H\) 上群结构的定义得到.

\medskip
\noindent\SGARefTarget{sga3:iv:remark:5:2:6}\textbf{注 5.2.6.}
在上述情形中, 若 \(f\) 的核恰好为 \(H\),
则 \(f\) 经由其分解的态射 \(G/H\to K\) 是单态射.
这由\SGARefLink{sga3:iv:proposition:3:3:4}{3.3.4} 立即得到.

对于群层, 可如下加强\SGARefLink{sga3:iv:proposition:4:4:11}{4.4.11}.

\medskip
\noindent\SGARefTarget{sga3:iv:proposition:5:2:7}\textbf{命题 5.2.7.}
\textit{设 \(G\) 为群层, \(H\) 为其正规子群层.
对 \(G\) 的每个包含 \(H\) 的子群层 \(K\),
令 \(K'=K/H\), 将其视为 \(G'=G/H\) 的子群. 则}
\[
  K=K'\times_{G'}G,
\]
\textit{且映射}
\[
  K\longmapsto K/H,
  \qquad
  K'\longmapsto K'\times_{G'}G
\]
\textit{给出 \(G\) 的包含 \(H\) 的子群层与 \(G'\) 的子群层之间的双射.
在此对应下, \(G\) 与 \(G'\) 的正规子群层相互对应.}

第一个断言容易由\SGARefLink{sga3:iv:proposition:4:4:11}{4.4.11}
和\SGARefLink{sga3:iv:lemma:3:2:4}{3.2.4} 得到.
只需证明 \(K\) 在 \(G\) 中正规当且仅当 \(K'\) 在 \(G'\) 中正规.
若 \(K\) 在 \(G\) 中正规, 则预层 \(i(K)/i(H)\)
在 \(i(G)/i(H)\) 中正规; 由通常的论证, 对伴随层也有同样结论.
反之, 若 \(K'\) 在 \(G'\) 中正规,
则纤维积 \(K'\times_{G'}G\) 在 \(G\) 中正规, 这显然成立.

现在设 \(L\) 为 \(G\) 的任意子群层.
用 \(\overline L\) 表示 \(L\) 对 \(H\) 定义的等价关系的饱和,
也记为 \(\overline L=L\cdot H\).

\medskip
\noindent\SGARefTarget{sga3:iv:proposition:5:2:8}\textbf{命题 5.2.8.}
\textit{在上述假设下, \(L\cdot H\) 是 \(G\) 的包含 \(H\) 的子群层,
且 \(L\) 在 \(G/H\) 中的像可等同于}
\[
  (L\cdot H)/H\simeq L/(H\cap L).
\]

% Source page: Exp4 component-local p. 60; combined-reader p. 248.
事实上, 用 \(L'\) 表示 \(L\) 在 \(G/H\) 中的像层.
在\SGARefLink{sga3:iv:proposition:5:2:7}{前一命题}的对应下,
它是 \(G/H\) 中对应于 \(L\cdot H\) 的子群层.
由于 \(L\to L'\) 是覆盖态射, 从而是层的普遍有效满态射,
由\SGARefLink{sga3:iv:proposition:4:4:9}{4.4.9} 可知,
\(L'\) 可等同于 \(L\) 对 \(L\to L'\) 的核的商,
而此核显然为 \(H\cap L\).

最后, 考虑群层 \(G\), 其子群层 \(K\), 以及 \(K\) 的正规子群层 \(H\).
首先在 \(G/H\) 上定义群层 \(H\backslash K\;(=K/H)\) 的右作用.
群 \(K\) 通过右平移作用于 \(G\).
由于 \(H\) 在 \(K\) 中正规, 此作用与 \(H\) 的作用所定义的等价关系相容,
从而诱导 \(K\) 在 \(G/H\) 上的作用, 即从 \(K\) 的反群出发的态射
\[
  K^{\circ}\longrightarrow\operatorname{Aut}(G/H)
\]
由于目标是层 (\SGARefLink{sga3:iv:corollary:4:5:13}{4.5.13}),
且此态射在 \(H\) 上平凡, 它经由 \(K/H\) 分解, 定义了所需作用.
\(G\) 在自身上的左、右作用可交换;
故 \(G\) 和 \(K/H\) 在 \(G/H\) 上的作用可交换.

\medskip
\noindent\SGARefTarget{sga3:iv:proposition:5:2:9}\textbf{命题 5.2.9.}
\textit{在上述假设下, \(K/H\) 自由地右作用于 \(G/H\),
且存在 \(G\)-对象的典范同构}
\[
  (G/H)/(K/H)\simeq G/K.
\]
\textit{若 \(K\) 在 \(G\) 中正规, 则由\SGARefLink{sga3:iv:proposition:5:2:7}{5.2.7},
\(K/H\) 在 \(G/H\) 中正规; 此时该同构保持两边的群结构.}

有预层的同构
\[
  \bigl(i(G)/i(H)\bigr)\big/\bigl(i(K)/i(H)\bigr)
  \xleftarrow{\sim} i(G)/i(K),
\]
它保持 \(i(G)\)-对象结构. 对此关系应用函子 \(a\), 即得所述结果.

\medskip
\noindent\SGARefTarget{sga3:iv:corollary:5:2:10}\textbf{推论 5.2.10.}
\textit{设 \(G\) 为 \(\mathscr C\)-群, \(K\) 为 \(G\) 的
\(\mathscr C\)-子群, \(H\) 为 \(K\) 的正规 \(\mathscr C\)-子群.
设 \((M)\) 为满足公理
\SGARefLink{sga3:iv:para:4:6:3:item:a:01}{(a)}--\SGARefLink{sga3:iv:para:4:6:3:item:e:mathcal:t:01}{\((e_{\mathcal T})\)}
的 \(\mathscr C\) 中的一族态射.
假设 \(H\) 在 \(G\) (相应地, \(K\)) 上的右作用是 \((M)\)-有效的.
则 \(K/H\) 自然且自由地右作用于 \(G/H\), 且此作用与 \(G\) 的作用可交换.
下列条件等价:}
\begin{enumerate}
  \item[(i)]\SGARefTarget{sga3:iv:corollary:5:2:10:item:i:01} \textit{\(K\) 在 \(G\) 上的作用是 \((M)\)-有效的;}
  \item[(ii)]\SGARefTarget{sga3:iv:corollary:5:2:10:item:ii:01} \textit{\(K/H\) 在 \(G/H\) 上的作用是 \((M)\)-有效的.}
\end{enumerate}
\textit{在这些条件下, 存在带有 \(G\) 作用的对象的同构\footnote{编者注: \(\mathscr C\) 的对象.} }
\[
  (G/H)/(K/H)\simeq G/K.
\]

% Source page: Exp4 component-local p. 61; combined-reader p. 249.
\SGARefTarget{sga3:iv:subsection:5:3}\subsection*{5.3. 有效性判据的应用: Noether 定理}
\addcontentsline{toc}{subsection}{\SGARefLink{sga3:iv:subsection:5:3}{5.3}. 有效性判据的应用: Noether 定理}
\footnote{编者注: 事实上, \SGARefLink{sga3:iv:proposition:5:2:8}{5.2.8}--\SGARefLink{sga3:iv:corollary:5:2:10}{5.2.10}
中建立的同构才称得上``Noether 同构定理''.}
设 \(\mathscr C\), \(\mathcal T\) 和 \((M)\) 如常.
设 \((N)\) 为满足\SGARefLink{sga3:iv:subsection:4:7}{4.7} 的公理
\SGARefLink{sga3:iv:subsection:4:7:item:a:01}{(a)} 和\SGARefLink{sga3:iv:subsection:4:7:item:f:m:01}{\((f_M)\)} 的态射族.
结合\SGARefLink{sga3:iv:proposition:5:2:7}{5.2.7} 与\SGARefLink{sga3:iv:proposition:4:7:2}{4.7.2},
得到\SGARefLink{sga3:iv:proposition:5:3:1}{下述结果}.

\medskip
\noindent\SGARefTarget{sga3:iv:proposition:5:3:1}\textbf{命题 5.3.1.}
\textit{设 \(G\) 为 \(\mathscr C\)-群,
\(H\) 为 \(G\) 的正规 \(\mathscr C\)-子群, \((M)\)-有效地作用于 \(G\).
对 \(G\) 的每个包含 \(H\) 且满足 \(K\to G\) 属于 \((N)\) 的
\(\mathscr C\)-子群 \(K\), \(H\) 在 \(K\) 上的作用是 \((M)\)-有效的,
且商 \(K/H=K'\) 是 \(G/H=G'\) 的 \(\mathscr C\)-子群,
满足 \(K'\to G'\) 属于 \((N)\).
映射 \(K\mapsto K'\) 是下列两个集合之间的双射:
\(G\) 的包含 \(H\) 且使 \(K\to G\) 属于 \((N)\) 的
\(\mathscr C\)-子群 \(K\) 所组成的集合,
以及 \(G'\) 的使 \(K'\to G'\) 属于 \((N)\) 的
\(\mathscr C\)-子群 \(K'\) 所组成的集合. 其逆为}
\[
  K'\longmapsto K\times_{G'}G.
\]
\textit{在此对应下, \(G\) 与 \(G'\) 的正规子群相互对应.}

\medskip
\noindent\SGARefTarget{sga3:iv:corollary:5:3:2}\textbf{推论 5.3.2.}
\textit{若 \(H\to G\) 属于 \((N)\), 则 \(\mathscr C\) 有终对象 \(e\),
且单位截面 \(e\to G/H\) 属于 \((N)\).}\footnote{编者注: 这由对 \(K=H\) 应用\SGARefLink{sga3:iv:proposition:5:3:1}{此命题}得到.}

% Translation coverage: complete Section 5 (source 460 lines).
% Audit: 395 math expressions, 43 SGARefTarget keys, 65 SGARefLink keys
% and 8 footnotes retained; math-expression multiset unchanged exactly.
